% ============================================================
%  MTA Transit Pulse — Project Report
%  Data Management Course
% ============================================================
\documentclass[12pt, letterpaper]{article}

% --- Packages -----------------------------------------------
\usepackage[utf8]{inputenc}
\usepackage[T1]{fontenc}
\usepackage{lmodern}
\usepackage[margin=1in]{geometry}
\usepackage{setspace}
\usepackage{titlesec}
\usepackage{titling}
\usepackage{fancyhdr}
\usepackage{hyperref}
\usepackage{graphicx}
\usepackage{float}
\usepackage{booktabs}
\usepackage{longtable}
\usepackage{listings}
\usepackage{xcolor}
\usepackage{caption}
\usepackage{subcaption}
\usepackage{amsmath}
\usepackage{parskip}
\usepackage[backend=biber, style=apa]{biblatex}
\addbibresource{references.bib}

% --- Colors -------------------------------------------------
\definecolor{subwayBlue}{RGB}{0, 57, 166}
\definecolor{codeGray}{RGB}{245, 245, 245}
\definecolor{codeGreen}{RGB}{0, 128, 0}
\definecolor{codeComment}{RGB}{100, 100, 100}
\definecolor{codePurple}{RGB}{128, 0, 128}

% --- Code listing style -------------------------------------
\lstdefinestyle{pythonstyle}{
    backgroundcolor=\color{codeGray},
    commentstyle=\color{codeComment}\itshape,
    keywordstyle=\color{subwayBlue}\bfseries,
    stringstyle=\color{codeGreen},
    basicstyle=\ttfamily\small,
    breakatwhitespace=false,
    breaklines=true,
    captionpos=b,
    keepspaces=true,
    numbers=left,
    numbersep=8pt,
    numberstyle=\tiny\color{codeComment},
    showspaces=false,
    showstringspaces=false,
    showtabs=false,
    tabsize=4,
    frame=single,
    rulecolor=\color{gray!30},
    language=Python
}

\lstdefinestyle{sqlstyle}{
    backgroundcolor=\color{codeGray},
    commentstyle=\color{codeComment}\itshape,
    keywordstyle=\color{subwayBlue}\bfseries,
    stringstyle=\color{codeGreen},
    basicstyle=\ttfamily\small,
    breaklines=true,
    captionpos=b,
    numbers=left,
    numbersep=8pt,
    numberstyle=\tiny\color{codeComment},
    frame=single,
    rulecolor=\color{gray!30},
    language=SQL
}

% --- Section formatting -------------------------------------
\titleformat{\section}
  {\large\bfseries\color{subwayBlue}}
  {\thesection}{1em}{}[\titlerule]

\titleformat{\subsection}
  {\normalsize\bfseries}
  {\thesubsection}{1em}{}

\titleformat{\subsubsection}
  {\normalsize\itshape}
  {\thesubsubsection}{1em}{}

% --- Header / Footer ----------------------------------------
\pagestyle{fancy}
\fancyhf{}
\lhead{\small MTA Transit Pulse}
\rhead{\small Data Management — \the\year}
\cfoot{\thepage}
\renewcommand{\headrulewidth}{0.4pt}

% --- Hyperlink style ----------------------------------------
\hypersetup{
    colorlinks=true,
    linkcolor=subwayBlue,
    urlcolor=subwayBlue,
    citecolor=subwayBlue,
    pdftitle={MTA Transit Pulse — Data Management Project},
}

% --- Title block --------------------------------------------
\title{
    \vspace{-1cm}
    {\large Data Management — Course Project Report}\\[0.4cm]
    \rule{\linewidth}{1.5pt}\\[0.3cm]
    {\LARGE\bfseries\color{subwayBlue} MTA Transit Pulse}\\[0.2cm]
    {\large A Data Pipeline for NYC Subway Ridership Analysis}\\
    \rule{\linewidth}{1pt}
}
\author{
    Your Name Here\\
    \small Student ID: XXXXXXX\\
    \small Department of Computer Science / Your Department\\
    \small Stony Brook University
}
\date{\today}

% ============================================================
\begin{document}
% ============================================================

\maketitle
\thispagestyle{empty}

% --- Abstract -----------------------------------------------
\begin{abstract}
\noindent
This report documents the design and implementation of \textit{MTA Transit Pulse},
a practical data management system built around New York City subway turnstile data
published by the Metropolitan Transportation Authority (MTA). The project demonstrates
a complete data lifecycle: collection from public sources, cleaning and validation,
relational storage in SQLite, SQL-driven analysis, automated backup, and interactive
visualization through a Streamlit dashboard. The system is version-controlled on GitHub
and documented throughout using Markdown and this report. Key findings include
ridership patterns across subway lines, peak usage hours, and station-level traffic
distributions.
\end{abstract}

\vspace{0.5cm}
\tableofcontents
\newpage

% ============================================================
\section{Introduction}
% ============================================================
%  Section: Introduction & Project Scope
% ============================================================

The New York City subway system is one of the largest and most complex urban
transit networks in the world, serving millions of riders daily across
\textbf{472 stations} and \textbf{27 lines}. Understanding ridership patterns,
identifying peak usage periods, and tracking station-level traffic are problems
with real operational significance for the MTA and for city planners.

This project, \textit{MTA Transit Pulse}, builds a complete data management pipeline
around the MTA's publicly available weekly turnstile data. Rather than simply
downloading and plotting a spreadsheet, the goal is to design and implement a
\textit{system}: one that collects, validates, stores, queries, and presents data
in a reproducible and maintainable way.

\subsection{Motivation}

% Write 2-3 sentences about why this project is interesting to you personally
% and why subway data is a meaningful dataset to work with.
\textit{[Add your personal motivation here.]}

\subsection{Project Objectives}

The system is designed to satisfy the following objectives:

\begin{enumerate}
    \item Automate the collection of MTA turnstile CSV data from the public portal.
    \item Clean and validate raw data, handling known MTA data quality issues.
    \item Store the processed data in a normalized SQLite relational database.
    \item Enable SQL-driven analysis of ridership patterns by line, station, and time.
    \item Provide automated database backups and version-controlled source code.
    \item Deliver an interactive Streamlit dashboard for non-technical data exploration.
    \item Document the full pipeline in a reproducible GitHub repository.
\end{enumerate}

\subsection{Scope and Limitations}

The project focuses on \textbf{entry and exit counts} from MTA turnstiles, which serve
as a proxy for ridership. The initial dataset covers
\textit{[insert your date range once you decide it]}.
Real-time train positions and delay data are outside the scope of this version,
though the database schema is designed to accommodate future enrichment.


% ============================================================
\section{Project Scope and Objectives}
\label{sec:scope}
% ============================================================
%  Section: Introduction & Project Scope
% ============================================================

The New York City subway system is one of the largest and most complex urban
transit networks in the world, serving millions of riders daily across
\textbf{472 stations} and \textbf{27 lines}. Understanding ridership patterns,
identifying peak usage periods, and tracking station-level traffic are problems
with real operational significance for the MTA and for city planners.

This project, \textit{MTA Transit Pulse}, builds a complete data management pipeline
around the MTA's publicly available weekly turnstile data. Rather than simply
downloading and plotting a spreadsheet, the goal is to design and implement a
\textit{system}: one that collects, validates, stores, queries, and presents data
in a reproducible and maintainable way.

\subsection{Motivation}

% Write 2-3 sentences about why this project is interesting to you personally
% and why subway data is a meaningful dataset to work with.
\textit{[Add your personal motivation here.]}

\subsection{Project Objectives}

The system is designed to satisfy the following objectives:

\begin{enumerate}
    \item Automate the collection of MTA turnstile CSV data from the public portal.
    \item Clean and validate raw data, handling known MTA data quality issues.
    \item Store the processed data in a normalized SQLite relational database.
    \item Enable SQL-driven analysis of ridership patterns by line, station, and time.
    \item Provide automated database backups and version-controlled source code.
    \item Deliver an interactive Streamlit dashboard for non-technical data exploration.
    \item Document the full pipeline in a reproducible GitHub repository.
\end{enumerate}

\subsection{Scope and Limitations}

The project focuses on \textbf{entry and exit counts} from MTA turnstiles, which serve
as a proxy for ridership. The initial dataset covers
\textit{[insert your date range once you decide it]}.
Real-time train positions and delay data are outside the scope of this version,
though the database schema is designed to accommodate future enrichment.

% NOTE: update this to reference a dedicated scope section
% once you separate them, or merge scope into introduction.

% ============================================================
\section{Data Collection}
\label{sec:data_collection}
% ============================================================
%  Section: Data Collection
% ============================================================

\subsection{Data Source}

The primary data source is the \textbf{MTA Turnstile Dataset}, published weekly
by the Metropolitan Transportation Authority at:
\url{http://web.mta.info/developers/turnstile.html}

Each weekly file is a CSV containing cumulative turnstile counter readings recorded
at approximately four-hour intervals across all subway stations. The columns in each
raw file are described in Table~\ref{tab:raw_schema}.

\begin{table}[H]
\centering
\caption{Raw MTA Turnstile CSV Column Schema}
\label{tab:raw_schema}
\begin{tabular}{@{}lll@{}}
\toprule
\textbf{Column} & \textbf{Type} & \textbf{Description} \\
\midrule
C/A        & String  & Control area (booth identifier) \\
UNIT       & String  & Remote unit ID \\
SCP        & String  & Subunit / channel / position \\
STATION    & String  & Station name \\
LINENAME   & String  & Lines serving this station \\
DIVISION   & String  & Operating division (BMT, IND, IRT) \\
DATE       & String  & Date of reading (MM/DD/YYYY) \\
TIME       & String  & Time of reading (HH:MM:SS) \\
DESC       & String  & Reading type (REGULAR or RECOVR AUD) \\
ENTRIES    & Integer & Cumulative entry counter \\
EXITS      & Integer & Cumulative exit counter \\
\bottomrule
\end{tabular}
\end{table}

\subsection{Collection Strategy}

% Describe your collect.py script here once you've written it.
% Topics to address:
%   - How URLs are constructed (the MTA uses a predictable date-based pattern)
%   - How many weeks of data you downloaded
%   - Where raw files are stored (data/raw/)
%   - How you avoided re-downloading files already on disk

\textit{[Complete this subsection after writing collect.py]}

\subsection{Known Data Quality Issues}

The MTA turnstile data is known to contain several quality issues that must be
addressed during preprocessing:

\begin{itemize}
    \item \textbf{Counter resets:} Turnstile counters periodically reset to zero,
          producing negative delta values when subtracted from the previous reading.
    \item \textbf{Anomalous spikes:} Occasionally a single reading shows millions
          of entries, far beyond any plausible passenger count.
    \item \textbf{Inconsistent station names:} The same physical station may appear
          under slightly different name strings across different files.
    \item \textbf{Duplicate rows:} Some audit entries (DESC = RECOVR AUD) overlap
          with regular readings.
\end{itemize}

These issues are addressed in Section~\ref{sec:cleaning}.


% ============================================================
\section{Data Cleaning and Preprocessing}
\label{sec:cleaning}
% Write this section after completing clean.py
% Topics to cover:
%   - What problems you found in the raw data
%   - How you handled missing values
%   - How you normalized station names
%   - Row counts before and after cleaning
\begin{quote}
    \textit{[Complete this section after finishing clean.py]}
\end{quote}

% ============================================================
\section{Database Design and Storage}
\label{sec:database}
% ============================================================
%  Section: Database Design and Storage
% ============================================================

\subsection{Database System}

The project uses \textbf{SQLite} as its relational database engine. SQLite is an
embedded, serverless database that stores the entire database as a single
\texttt{.db} file. This choice was made for several practical reasons:

\begin{itemize}
    \item No server installation or configuration required.
    \item Native Python support via the built-in \texttt{sqlite3} module.
    \item Sufficient performance for analytical queries on millions of rows.
    \item Portability: the database file can be shared or recreated from the pipeline.
\end{itemize}

\subsection{Schema Design}

% Add your ER diagram here once you have it.
% Placeholder:
\begin{figure}[H]
    \centering
    % \includegraphics[width=0.85\textwidth]{figures/er_diagram.png}
    \fbox{\parbox{0.85\textwidth}{\centering\vspace{2cm}
    \textit{[Insert ER diagram here]}
    \vspace{2cm}}}
    \caption{Entity-Relationship Diagram for the MTA Transit Pulse database}
    \label{fig:er_diagram}
\end{figure}

The database consists of the following tables:

\subsubsection{\texttt{stations}}
Stores one row per unique station, normalized from the raw STATION column.

\subsubsection{\texttt{turnstiles}}
Stores one row per unique turnstile unit (identified by C/A + UNIT + SCP).
Each turnstile is linked to a station via foreign key.

\subsubsection{\texttt{readings}}
Stores the individual counter readings with parsed timestamps.
This is the largest table, with one row per four-hour interval per turnstile.

\subsubsection{\texttt{daily\_ridership}}
A derived/aggregated table that pre-computes daily net entries and exits per
station. This table is rebuilt by the analysis pipeline and powers the
Streamlit dashboard queries.

\subsection{Key SQL Patterns}

% Once you've written your queries, paste 1-2 representative ones here.
% Example placeholder:

\begin{lstlisting}[style=sqlstyle, caption={Example: Top 10 stations by average daily entries}]
SELECT
    s.station_name,
    AVG(dr.net_entries) AS avg_daily_entries
FROM daily_ridership dr
JOIN stations s ON dr.station_id = s.id
GROUP BY s.station_name
ORDER BY avg_daily_entries DESC
LIMIT 10;
\end{lstlisting}


% ============================================================
\section{Data Analysis}
\label{sec:analysis}
% ============================================================
%  Section: Analysis
% ============================================================

\subsection{Analysis Questions}

The analysis notebook addresses the following questions using SQL queries
executed from Jupyter via the \texttt{sqlite3} Python module:

\begin{enumerate}
    \item Which 10 stations have the highest average daily ridership?
    \item What hours of the day see the most subway entries system-wide?
    \item How does ridership vary by day of the week?
    \item Which subway lines carry the most passengers?
    \item Are there observable ridership drops on public holidays?
\end{enumerate}

% Complete each subsection below as you answer each question in your notebook.

\subsection{Station-Level Ridership}
\textit{[Insert findings and a figure here after running your notebook]}

\subsection{Temporal Patterns}
\textit{[Insert findings and a figure here after running your notebook]}

\subsection{Line-Level Comparison}
\textit{[Insert findings and a figure here after running your notebook]}


% ============================================================
\section{Backup and Security Measures}
\label{sec:backup}
% Topics to cover:
%   - Your backup script and schedule
%   - What is .gitignored and why
%   - How you handled the .env file
\begin{quote}
    \textit{[Complete this section after finishing your backup script]}
\end{quote}

% ============================================================
\section{Data Maintenance and Updates}
\label{sec:maintenance}
% Topics to cover:
%   - How new weekly turnstile files get ingested
%   - Deduplication strategy
%   - How you'd schedule automated updates
\begin{quote}
    \textit{[Complete this section after your pipeline is running end-to-end]}
\end{quote}

% ============================================================
\section{Version Control and Documentation}
\label{sec:versioning}
% Topics to cover:
%   - Git workflow you followed (commits, branches if any)
%   - README structure
%   - How this report relates to the codebase
\begin{quote}
    \textit{[Complete this section near the end of the project]}
\end{quote}

% ============================================================
\section{Visualization}
\label{sec:visualization}
\input{sections/visualization}

% ============================================================
\section{Streamlit User Interface}
\label{sec:ui}
% Topics to cover:
%   - How the app is structured
%   - What queries it runs against SQLite
%   - Screenshots of the interface
%   - How to run it locally
\begin{quote}
    \textit{[Complete this section after finishing streamlit\_app.py]}
\end{quote}

% ============================================================
\section{Conclusion}
\label{sec:conclusion}
\input{sections/conclusion}

% ============================================================
\newpage
\printbibliography[title={References}]

\end{document}
